\documentclass[12pt,a4paper]{article}

\usepackage[italian]{babel}  % Per la sillabazione e le intestazioni in italiano
\usepackage[utf8]{inputenc} % Encoding caratteri (accenti, ecc.)
\usepackage[T1]{fontenc}    % Font vector reali per caratteri accentati
\usepackage{hyperref}       % Per link e metadati del PDF

\title{
    \large\bfseries Relazione progetto per il corso:  \\[0.6em]
    \large Introduzione alla programmazione di smart contract per Ethereum \\[0.9em]
    \LARGE Marketplace per NFT nell'ambito di una filiera farmaceutica
}

\author{
    \textbf{Ariana Baci} \\[0.8em]
    \small Anno Accademico 2025/2026 \\[0.8em]
    \small Corso di Laurea in Informatica \\
    \small Università degli Studi di Verona
}

\date{\today}

\begin{document}

\maketitle
\newpage
\section{  \Large Interpretazione della traccia }
Il progetto nasce dall'idea di sfruttare la tecnologia blockchain per garantire la tracciabilità e l'autenticità dei prodotti. In particolare, in questo caso ho pensato al caso della filiera farmaceutica, dove può essere estremamente importante avere la garanzia che un
determinato prodotto derivi da uno specifico produttore.
La situazione che ho rappresentato è una semplificazione di un'ipotetica situazione reale.
Ci sono due contratti: uno appartiene al laboratorio che produce i farmaci, il secondo
si occupa di caricare gli annunci relativi. Idealmente quindi, il laboratorio, quando un farmaco viene reso disponibile, ne crea l'NFT corrispondente
e ne carica l'annuncio sul marketplace scelto per mettere in vendita i propri prodotti.
Il marketplace deve sia essere autorizzato dal contratto che crea gli NFT, sia autorizzare lui stesso quel contratto a pubblicare annunci. In questo caso
ho reso possibile per semplicità l'autorizzazione di un solo contratto, ma si potrebbe estendere per fare in modo che il marketplace gestisca 
una lista di vari produttori di farmaci. \newline Questo sistema risulta molto più sicuro rispetto 
alla semplice registrazione dei dati in un database centralizzato: grazie all'unicità degli NFT non c'è il rischio che vengano ad esempio riciclati codici identificativi per vendere prodotti contraffatti, mentre la decentralizzazione garantisce
che il produttore non modifichi i dati immutabili a posteriori e che i dati non smettano di essere fruibili se cade un server.
\section{  \Large Architettura del codice}
Il progetto si compone di due contratti: DrugNFT.sol crea gli NFT che rappresentano i farmaci
e DrugMarkenplace.sol si occupa di gestirne gli annunci di vendita.
\subsection{DrugNFT}
DrugNFT è il contratto che permette di mintare i farmaci.
Implementa lo standard ERC-721 (esteso tramite \texttt{ERC721URIStorage} di OpenZeppelin) per la gestione dei metadati e Ownable per il controllo degli accessi basato su un unico proprietario (in questo caso direttamente collegato al laboratorio
produttore). Grazie allo struct apposito e al mapping associa univocamente a ogni tokenId nome, numero di lotto, data di produzione, data di scadenza, produttore, acquirente e gli indicatori di stato (forSale, sold).
L'operazione di minting è riservata al proprietario del contratto (onlyOwner). Al suo interno, la funzione mintDrugNFT esegue controlli sulle date, 
aumenta il valore del tokenID prima di chiamare _safeMint per assicurare che non vi siano duplicati, emette l'evento DrugMinted così da notificare la creazione del nuovo NFT.
La funzione setMarketplace permette al produttore (onlyOwner) di registrare il contratto DrugMarketplace. Applica automaticamente l'approvazione globale (setApprovalForAll), consentendo al marketplace di trasferire il token al momento dell'acquisto senza richiedere approvazioni manuali ad ogni vendita.
Le funzioni markAsSold e markAsForSale sono definite onlyMarketplace (modificatore che approva l'esecuzione del corpo della funzione solo al marketplace) 
e permettono di variare i campi del NFT relativi allo stato di vendita in maniera istantanea e automatica.
Le funzioni getDrug, isExpired, isSold, isForSale permettono anche a contratti terzi di verificare varie informazioni sul prodotto.

\subsection{DrugMarketplace}
DrugMarketplace è il contratto che gestisce la compravendita dei farmaci.
Utilizza ReentrancyGuard per proteggere da attacchi di reentrancy.
Le funzioni setDrugNFT e unSetDrugNFT permettono all'amministratore (onlyOwner) di associare o disassociare il contratto DrugNFT di riferimento (che in questo caso
è unico ma realisticamente un marketplace potrebbe gestirne di più).
ListDrug consente solo al legittimo possessore dell'NFT di mettere in vendita il farmaco impostando un prezzo (price>0).
Verifica che il marketplace sia stato accreditato da DrugNFT, aggiorna lo stato interno (isActive = true), sincronizza il flag su DrugNFT (markAsForSale) ed emette l'evento DrugListed.
BuyDrug permette a qualsiasi indirizzo (diverso dal venditore) di acquistare un annuncio attivo e non scaduto;
trasferisce la proprietà dell'NFT dal venditore all'acquirente via safeTransferFrom, aggiorna i flag di vendita (markAsSold, markAsForSale = false) e disattiva l'annuncio,
nvia il pagamento al venditore tramite call e rimborsa automaticamente all'acquirente l'eventuale ETH inviato in eccesso (msg.value>price).
CancelListing e updatePrice permettono solo al venditore di annullare un annuncio attivo (es. per imminente scadenza o ritiro dal mercato) o di aggiornarne il prezzo di vendita.
GetListing, getPrice, getSeller permettono consultare rapidamente i dettagli degli annunci.

\section{Test }
I test sono specifici per i due contratti. Quelli in typescript sono nella cartella "test" mentre quelli in solidity si trovano in "contracts".
\subsection{Test in Typescript}
\subsubsection[short]{DrugNFT}
Inizializzazione del Contratto: Verifica che il nome ("DrugTraceability") e il simbolo ("DRUG") siano impostato correttamente.
Conferma che il ruolo di owner appartenga al wallet del producer.
Creazione/Minting (mintDrugNFT):
Caso di successo: Verifica che il producer riesca a coniare un NFT e che la struct Drug salvata contenga i valori corretti.
Verifica che l'account other non autorizzato fallisca se tenta di mintare un NFT.
Verifica che la transazione fallisca se expirationDate≤productionDate oppure se il farmaco ha una data di scadenza nel passato.

Interazione con il Marketplace:
verifica che solo l'owner del contratto possa invocare setMarketplace.
e che solo il marketplace autorizzato possa invocare markAsSold e markAsForSale.
\subsubsection[short]{DrugMarketplace}
Sfrutta networkHelpers creando una fixiture (deployMarketplaceFixture) così da partire sempre dallo stesso stato senza dover rieseguire ogni volta i vari deploy dei contratti.
Verifche svolte: 
\begin{itemize}
    \item setDrugNFT / unSetDrugNFT: verifica che solo l'amministratore possa impostare o azzerare (address(0)) l'indirizzo del contratto DrugNFT.
    \item listDrug: verifica che la pubblicazione di un annuncio sia permessa solo al proprietario dell'NFT.
    \item buyDrug. Si verifica l'impossibilità di: acquistare annunci revocati/inattivi, acquistare un farmaco di cui si è proprietari, acquistare con fondi insuffiienti o farmaci scaduti. In caso di successo invece, si verifica che un acquisto valido trasferisca il token, aggiorni lo stato a isSold = true ed esegua la transazione.
    \item updatePrice, getListing, etc: controlla che solo il venditore possa modificare il prezzo con updatePrice e 
 la corretta estrazione dei dati da parte delle funzioni getPrice, getSeller e getListing.
    \item 
\end{itemize}
\subsection{Test in Solidity}
\subsubsection[short]{DrugNFT}
\begin{itemize}
\item Configurazione (setUp): Istanzia DrugNFT impostando l'indirizzo fittizio producer (0x1111) come proprietario iniziale.
\item Minting : simula la chiamata da parte del produttore autorizzato sfruttando vm.prank(producer), crea un NFT,
verifica che il tokenId generato sia 0 (primo token mintato) e che le variabili interne alla struct Drug (name, lotNumber, producer, sold = false) corrispondano esattamente ai dati inseriti.
\item Scadenze (test_IsExpired): sfruttando vm.warp simula l'avanzare del tempo dell'EVM, facendo trascorrere 101 secondi in modo da far scadere il farmaco appositamente mintato con scadenza a 50 secondi.
Verifica che la funzione isExpired ritorni true a seguito del superamento della data di scadenza.
\end{itemize}

\subsubsection[short]{DrugMarketplace}
\begin{itemize}
\item Configurazione (setUp): Esegue il deploy di entrambi i contratti (DrugNFT e DrugMarketplace). Stabilisce i riferimenti incrociati (setMarketplace su DrugNFT e setDrugNFT su DrugMarketplace) usando vm.prank(producer).
\item Test flusso di vendita (testVendita): usando vm.startPrank(producer), il produttore conia il token #0 e crea l'annuncio su DrugMarketplace per 100 wei.
Con vm.deal(buyer, 1 ether) aggiunge i fondi necessari all'acquirente, esegue poi l'acquisto con marketplace.buyDrug sotto vm.startPrank(buyer).
\item Accerta il trasferimento di proprietà dell'NFT in favore dell'acquirente (drugNFT.ownerOf(tokenId) == buyer), estrae la tupla da marketplace.getListing tramite la sintassi a parentesi tonde () e verifica che isActive e isForSale siano resettati a false.
\end{itemize}

\end{document}